\chapter{Reflection and AI use}
\label{chap:project-management}

\section{Challenges and trade-offs}
With limited time, a key challenge was deciding where testing effort
would be most useful. Priorities were set early so that time could be
allocated to discount and delivery rules, their interactions, and the
validation paths that could affect an order.

\begin{itemize}
    \item Representative input combinations and boundary cases were
    selected in place of an exhaustive cross-product. This kept the
    tests focused on pricing decisions, while leaving some combinations
    untested.
    \item The structural pass prioritised stock checks and invalid
    cart and delivery inputs. Additional tests for customer and product
    setters were deferred because they offered less evidence about
    the amount payable.
\end{itemize}

With more time, the decision-to-decision sketches for selected methods
would have been extended into detailed DD graphs for the remaining
methods. The tests could then have been mapped to individual edges to
identify further reachable paths.

\section{AI assistance and validation}
\subsection{Purpose, scope, and prompts}
% Explain why you used AI and what it helped with. Include the actual prompts
% required by the brief. Record help with the LaTeX setup and evidence table.

\subsection{Independent validation}
% Explain how you checked AI-assisted scenarios, expected values, assertions,
% and results against the brief. Describe checks you actually carried out.

